\section{\unclass Test Structure}  

\subsection{\unclass Tests Description}  

\guidance{

\unclass TIMELINE: Should be started by PDR and finalized at CDR.

\unclass This section should generally not be modified.  However, if you have a good reason to deviate from this standard, your model may modify this format to better suite your testing needs with approval from, at a minimum; the model lead, the NGIC M\&S TPOC, and the analyst for the system.  
}
%%%%% BEGIN USER CONTENT %%%%%

\classlabelshort The tests for the \modelname{} model are split into a number of test families which represent multiple geometries or over-arching test conditions, and these test families contain test categories that demonstrate multiple requirements. Miscellaneous test categories that include tests that don’t fulfill any requirements but show useful results, are also included. Each family is split into a series of tests, which will fulfill all of the requirements for the corresponding requirement set.

\unclass System tests assume that all unit tests pass for each component used in the system test.

\unclass Fully completed tests will be marked with a \checkmark, and their results will be explored in more detail below.  

Each Test Case describes the \gls{kpi}s that test assesses.  A table of all \gls{kpi}s and the associated equations appears in the \gls{uam} as Table \#.

\subsection{\classlabelshort Test Numbering Scheme}  


\guidance{	\unclass TIMELINE: Should be started by PDR and finalized at CDR.\\
	
\unclass Briefly state the schema used to number and delineate test cases.
}
\unclass System tests will be named according to a set of major Test Families which will be described below using the designation "TF1" to mean "Test Family 1". Some requirements will take more than a single test to meet, and some tests can be designed to satisfy more than one requirement, so a 1-1 tests-requirements match-up will not exist.  Please refer to the requirements traceability matrix in section \ref{RTM} for more details on tests to requirements mapping for model requirements. 




